
\chapter{Introduction}


\ifpdf
    \graphicspath{{1_introduction/figures/PNG/}{1_introduction/figures/PDF/}{1_introduction/figures/}}
\else
    \graphicspath{{1_introduction/figures/EPS/}{1_introduction/figures/}}
\fi


\section{Overview}


{\color{red} 
This Final Year Project compares two families of deep reinforcement learning algorithms: value-based methods, which learn by estimating the value of situations, and actor-critic methods, which learn by directly improving the strategy that works. Three experiments of increasing scope and depth test both families under identical conditions --- same framework, same default settings, no per-algorithm tuning --- so that differences in results reflect genuine algorithmic character rather than configuration effort.

The results show that actor-critic methods are the more reliable default: they produce strong, consistent results across all tested settings with no special setup. Value-based methods are not weaker --- they are more demanding. They require more training time, more careful configuration, and greater domain knowledge to reach their potential. This distinction matters for practice: the right question is not which family is better, but which is appropriate given the available budget and expertise.

A secondary finding is methodological: from the same data, different evaluation metrics can produce opposite conclusions about which algorithm family wins. This project demonstrates that a single aggregate number is insufficient --- and potentially misleading --- when comparing reinforcement learning algorithms.
}

%%%%%%%%%%%%%%%%%%%%%%%%%%%%%%%%%%%%%%%%%%%%%%%%%%%%%
% \textbf{2nd and 3rd paragraphs introduce RL} \\
%%%%%%%%%%%%%%%%%%%%%%%%%%%%%%%%%%%%%%%%%%%%%%%%%%%%%

%[P2: Introduce RL]
\textbf{Reinforcement learning} is learning what to do — how to map situations to actions — so as to maximize a numerical reward signal. The learner is not told which actions to take, but instead must discover which actions yield the most reward by trying them. In the most interesting and challenging cases, actions may affect not only the immediate reward but also the next situation and, through that, all subsequent rewards. These two characteristics—trial-and-error search and delayed reward—are the two most important distinguishing features of reinforcement learning. \cite{sutton}

%[P3: RL History, Perceptron Learning, Associative Learning] --> MOVE TO CHAPTER 2
%Reinforcement learning originates from two major sources: animal learning and optimal control. Early research in the 20th century focused on understanding how animals adapt behavior through trial-and-error processes.
%Edward Thorndike’s experiments with cats in 1911 established the principles of behavioral reinforcement, while Pavlov’s work in 1927 laid the groundwork for stimulus-response associations. Skinner’s operant conditioning (1938) extended these ideas, demonstrating how behavior is shaped through external reinforcements.
%By 1949, Donald Hebb introduced the concept of synaptic strengthening, influencing modern neural networks. Finally, the Rescorla-Wagner model (1972) formalized learning dynamics, providing a predictive framework for associative learning. \cite{rlhistoryPod}

%[P3: RL continued]
Learning how to play chess can be a good example of Reinforcement learning, by making discrete steps in the environment, which are called actions (moving a piece) and getting back a reward (how good your current position), we can determine how good our move was, and vice versa, learn to make better moves. The difference from other methods unique to RL is that we give a score to the actions taken rather than saying correct/incorrect as is done in supervised learning. In RL, you need to balance between exploitation of already learned actions with good reward and exploration to potentially find even better actions to exploit for better reward in the long-term goal.

%[P4: RL continued part e]
Reinforcement Learning is designed to solve problems that can
be modeled as sequential decision processes. \cite{barto1989learning}
Learning to play chess is a sequential decision process, as every action taken by the player (agent) transitions the environment (chess board) to a new state. For example, moving the knight to the right cell may lead to checkmate of the enemy's king and winning or open an opportunity for your opponent to attack you instead, which may lead to losing the game.


%%%%%%%%%%%%%%%%%%%%%%%%%%%%%%%%%%%%%%%%%%%%%%%%%%%%%
% \textbf{4th and 5th paragraphs introduce CNNs and DQN} \\
%%%%%%%%%%%%%%%%%%%%%%%%%%%%%%%%%%%%%%%%%%%%%%%%%%%%%

%[P5: Search Spaces]
The search space for the game of Connect 4 is $10^{12}$ and for chess is $10^{40}$. It is difficult to enumerate each board in such large search spaces using brute force approach such as breadth first or depth first search algorithms. But this was achieved  by IBM’s Deep Blue super computer beat the chess world champion Garry Kasparov back in 1996. However, the search space for other games such as Go is significantly larger and brute search approaches do not work. Instead one has to use function approximators and deep learning architectures in particular given the recent success of Convolutional Neural Networks (CNNs) in computer vision. 

%it wasn’t until fairly recently that computers were able to reliably perform seemingly trivial tasks such as detecting a puppy in a picture or recognizing spoken words. Why are these tasks so effortless to us humans? The answer lies in the fact that perception largely takes place outside the realm of our consciousness, within specialized visual, auditory, and other sensory modules in our brains. By the time sensory information reaches our consciousness, it is already adorned with high-level features; \textbf{for example, when you look at a picture of a cute puppy, you cannot choose not to see the puppy, not to notice its cuteness. Nor can you explain how you recognize a cute puppy;} it’s just obvious to you. Thus, we cannot trust our subjective experience: perception is not trivial at all, and to understand it we must look at how our sensory modules work. \cite{geron}

%Convolutional neural networks (CNNs) emerged from the study of the brain’s visual cortex, and they have been used in computer image recognition since the 1980s.

%[P6: CNNs]
Convolutional Neural Networks (CNNs) are a type of neural network designed specifically for data that comes in grid-like formats, such as images. Instead of treating every input value independently, a CNN works with small local regions of the image and learns filters that can detect useful patterns. These filters slide across the input, allowing the network to find the same pattern in different locations. This idea of local connections and shared weights is one of the core principles of CNNs, along with pooling and stacking multiple layers to build deeper feature representations \cite{lecun:hal-04206682}

%[P7: Value Functions and DQN]
There are two major families of algorithms in the RL family tree. One branch is Value Functions which estimate the value of the likely long term reward from the current position or from slecting an action in the current position. Typical of thiese is Q Learning Method, and its implementation on deep learning architecture called Deep Q Network (DQN). This was used by DeepMind with layers of tiled convolutional filters to mimic the effects of receptive fields. Reinforcement learning is unstable or divergent when a nonlinear function approximator such as a neural network is used to represent Q. This instability comes from the correlations present in the sequence of observations, the fact that small updates to Q may significantly change the policy of the agent and the data distribution, and the correlations between Q and the target values. The method can be used for stochastic search in various domains and applications.\cite{li2023reinforcement} \cite{matzliach2022detection} \\

% \underline{Paragraph need on POLICY GRADIENT APPROACHES}

The other major branch uses the Policy Gradient approach. Instead of learning value functions, this group of algorithms parameterizes the policy directly and returns probabilities for each possible action in the current state. Along with other distinguishing features that will be discussed later, Policy Gradient methods allow us to train networks in spaces which are not handled well by value-based algorithms—such as stochastic environments or continuous action spaces, where there is no certainty of which action is best, where functions like argmax do not work, or where the number of possible actions is enormous.

% \underline{Paragraph needed on gym}

For Prototypes developed in this Report I used Gymnasium which is a maintained fork of OpenAI Gym and is now the recommended standard for developing and evaluating reinforcement learning agents. Gymnasium is described as "an API standard for reinforcement learning with a diverse collection of reference environments" and provides a simple, pythonic interface capable of representing general reinforcement learning problems. 
%%%%%%%%%%%%%%%%%%%%%%%%%%%%%%%%%%%%%%%%%%%%%%%%%%%%%


\section{Research Questions}

\textbf{RQ:} How does the performance of value-based Deep Q-Network agents compare to actor-critic agents (specifically PPO and A2C) in discrete action space environments, including their sample efficiency and training consistency during training?


These algorithms will be compared based on several criteria: 
\textbf{convergence speed} (number of episodes to reach the threshold)
\textbf{sample efficiency} (performance per interaction)
\textbf{training stability }(reward variance). 

These comparison criteria will serve as the basis for evaluating the trade-off between the fast, off-policy learning of value-based methods versus the stable, on-policy learning of actor-critic methods. \\

\textbf{Secondary Objective:} Achieve a solid understanding of the foundations of artificial intelligence, machine learning, and deep reinforcement learning.










% Steps
% \begin{enumerate}
%     \item Approximate a Systematic Literature Review
%     \item \subsection{Experimental Design}
%     \item 
%     \item Benchmarking and Visualing
% \end{enumerate}z





{\color{red}
this chpater needs check/update

\section{Methodology}
}
This section outlines the methodologies followed to ensure the successful
execution of this project.





\subsection{Systematic Literature Review}
To ensure the prototype development and experimentation tasks were grounded in existing research, a systematic literature review approximation of artificial intelligence, machine learning, reinforcement learning, value-based, and policy based reinforcement learning relevant texts and research papers was first conducted. Google Scholar, arXiv, and IEEE Xplore were electronic databases employed during the search process. Search terms used were "Artificial Intelligence", "Machine Learning", "Reinforcement Learning", "Deep Reinforcement Learning", "Deep Q-Network", "Actor-Critic", "Policy Gradient", "A2C", and "PPO". The inclusion criteria were articles and research papers that discussed at least one fundamental machine learning or reinforcement learning concept, or value-based or actor-critic reinforcement learning algorithms examined in this final year project. Educational video content and other online resources, such as framework documentation, were also queried where applicable. "Reinforcement Learning: An Introduction" by Andrew Barto and Richard S. Sutton, and "Hands-On Machine Learning with Scikit-Learn, Keras, and TensorFlow" by Aurélien Géron were two texts heavily relied upon when gaining knowledge of machine learning and reinforcement learning fundamentals.

The systematic literature review approximation can be broken down into five topics, as outlined below:
\begin{enumerate}

    \item The Foundations of Artificial Intelligence and Machine Learning: History, Semantics, PSSH, POE Model, Machine Learning Approaches.

    \item The Fundamental Principles and Approaches of Reinforcement Learning: Terminology, Exploitation vs Exploration, The Credit Assignment Problem, MDPs, The Bellman Equation, Dynamic Programming, GPI, Monte Carlo Algorithm, Temporal Difference Learning, Q-learning.

    \item Deep Learning: Threshold Logic Unit, Perceptron, Multilayer Perceptron, Backpropagation, Hyperparameters, Gradient Instability, Batch Normalization, Convolutional Neural Networks.

    \item Value-Based and Actor-Critic Reinforcement Learning Algorithms: Rationale, Deep Q-learning, Policy Gradients, Actor-Critic Methods, Advantage Estimation, On-Policy vs Off-Policy Learning, Training Stability.

    \item Atari Breakout and Statistical Significance: OpenAI, History, Rationale.

\end{enumerate}



\subsection{Experimental Design}
\label{subsec:experimental-design}

{\color{red} While studies comparing these algorithm families exist \citep{defuente2024comparative}, they typically involve per-algorithm hyperparameter tuning, making it difficult to separate algorithmic capability from configuration effort. This project takes a different approach: Stable-Baselines3 default hyperparameters are used for every algorithm across all three experiments, with no tuning of any kind. The importance of this distinction is illustrated by \citet{mnih2015human}, who acknowledge that their published DQN results relied on hyperparameters ``selected by performing an informal search,'' noting that ``it is conceivable that even better results could be obtained by systematically tuning the hyperparameter values.'' All results in this thesis reflect what each algorithm delivers out of the box --- an intentional choice to measure inherent algorithmic behaviour rather than tuning effort.} Moreover, implementing both algorithm families facilitated the validation of existing research and the expansion on recent findings. Establishing a basis for future research and fostering transparency were other motivating factors for prototype development.

In addition to discussing the main implementation components that comprise the DQN and Actor-Critic architectures, Chapter 3 provides a comprehensive overview of the high-level architectural decisions, equipment employed, and problems encountered during experimentation. Chapter 3 can be broken down into four topics:

\begin{enumerate}
\item High Level Architecture Decisions: Framework, Program Structure, Implementation Medium.
\item Code Implementation: Atari Breakout Environment, DQN and Actor-Critic Architectures, Recording Training Progress.
\item Equipment Employed.
\item Problems Encountered: GPU Driver Compatibility, Memory Leak.
\end{enumerate}




\subsection{Data Collection}

{\color{red} 
Data collection followed the same protocol across all three experiments. During training, performance was measured via periodic frozen-policy evaluations rather than live training rewards, since training rewards are algorithm-dependent and not a stable basis for comparison. At fixed step intervals, each agent was evaluated for a fixed number of episodes in a separate environment with no exploration noise, and the mean episodic return was recorded.

Across all experiments, training used Stable-Baselines3 (Section~\ref{sec:tooling}) with its built-in \texttt{EvalCallback}, which handles evaluation scheduling and logging automatically. Final performance was assessed with a dedicated post-training evaluation of 50 episodes per seed. All seeds were derived deterministically from a single master seed per experiment using NumPy's \texttt{SeedSequence}, ensuring reproducibility. Full details of evaluation frequency, episode counts, and normalisation are defined in Section~\ref{sec:eval-framework}.
}













\subsection{Data Analysis}

{\color{red} 
Results were analysed using the \texttt{rliable} statistical evaluation framework \citep{agarwal2021precipice} (Section~\ref{sec:eval-framework}). For each experiment, per-environment learning curves show performance over the training budget across all seeds, giving a picture of both speed and consistency. Final performance is summarised using interquartile mean (IQM) and pairwise probability of improvement, each reported with 95\% stratified bootstrap confidence intervals. Cross-environment aggregates are treated as descriptive summaries rather than definitive rankings, given the small number of environments per experiment. Full metric definitions and the normalisation procedure are provided in Section~\ref{sec:eval-framework}.
}



\subsection{Contributions}

{\color{red} 
This project makes three contributions to the empirical study of deep reinforcement learning.

\begin{enumerate}
\item A three-experiment benchmark comparing value-based and actor-critic algorithm families, using a single shared framework with default hyperparameters throughout (see Section~\ref{sec:tooling}). Over 80 training runs were conducted across classic-control and Atari environments, with results reproducible and consistent across seeds.

\item Application of the \texttt{rliable} statistical evaluation framework (Section~\ref{sec:eval-framework}) to a student-scale RL study, reporting interquartile mean, probability of improvement, and bootstrap confidence intervals rather than mean scores alone.

\item An observation that cross-environment aggregation metric choice affects the conclusion of the comparison: the same data produces different aggregate rankings under IQM versus pairwise probability of improvement. This is not a flaw in the methodology but a finding about the limits of single-number RL benchmarks.
\end{enumerate}
}





















\section{Project Plan}

The first semester was spent building the theoretical foundation. This involved working through the literature on reinforcement learning fundamentals --- MDPs, Bellman equations, Q-learning, policy gradients, and actor-critic methods --- as well as implementing several small-scale environments (Cliff Walking, Windy Gridworld, Random Walk) to build intuition before tackling anything larger. The semester ended with a working DQN agent on Atari Breakout as a proof of concept.

The second semester moved into the main comparison. Three experiments were run in sequence, each extending the scope of the previous one. Experiment~1 covered five algorithms on classic-control environments using MLP policies. Experiment~2 narrowed the focus to DQN vs PPO on Atari with CNN policies at a 5M-step budget. Experiment~3 extended that to 20M steps and added Seaquest as a third environment. The results of all three are reported in Chapters~4, 5, and 6 respectively.


\section{Overview of the Report}

{\color{red} 
Chapter~2 reviews the theoretical foundations of reinforcement learning --- the two algorithm families, the key algorithms from each, and the tools and evaluation framework used throughout. Chapter~3 presents the prototype work from the first semester: a DQN implementation on Atari Breakout that established the technical foundation for the experiments. The project GitHub repository contains the full source code for all experiments, commit history, and development statistics, and is linked in Chapter~3.

The three experiment chapters were not designed upfront as a fixed plan --- each arose from the findings and open questions of the previous one. Chapter~4 evaluates five algorithms across both families on classic-control environments. Algorithm and environment descriptions are provided in Chapter~\ref{chap:lit-review}. Chapter~5 narrows the comparison to one representative per family and scales to pixel-based Atari games. Chapter~6 extends that comparison to a longer training budget and a wider set of environments. Each chapter follows the same structure: setup, hypotheses, results, key findings, threats to validity, and discussion.

Chapter~7 synthesises the findings across all three experiments and provides a direct answer to the research question. Chapter~8 closes the report: it restates the conclusions plainly, reflects on what made this project difficult, and outlines directions for future work.
}






































